% Standard pen-and-paper worksheet preamble, in the form m01 and m02 use.
% Compile with: xelatex -interaction=nonstopmode exercise.tex   (run twice)
%
% 14pt, not 12pt: the sheet is read with a pencil in hand. There are no blanks
% and no ruled answer gaps anywhere on it -- questions are asked in words and
% answered on the back -- so the paper those would have cost goes to the text
% and to the drawings.
\documentclass[a4paper, 14pt]{extarticle}
\usepackage[top=0.75in, bottom=0.7in, left=0.85in, right=0.85in]{geometry}
\usepackage{amsmath}
\usepackage{amssymb}
\usepackage{array}
\usepackage{graphicx}
\usepackage[hidelinks]{hyperref}   % hidelinks, or the print carries colour boxes
\usepackage{tcolorbox}
\usepackage{enumitem}
\usepackage{etoolbox}
\usepackage{tikz}
\usepackage{fontspec}
\usetikzlibrary{calc,decorations,patterns,arrows,decorations.pathmorphing,positioning}
\definecolor{pltblue}{HTML}{1F77B4}

% Body face: Charter, and the four ways to get it. XCharter *is* Charter and
% ships with a full TeX Live, so it is asked for first and by file -- a CI
% runner answers yes to \IfFontExistsTF{Charter}, then cannot typeset it, which
% is how these sheets stopped building there.
\IfFontExistsTF{XCharter-Roman.otf}{
  \setmainfont{XCharter}[
    Extension      = .otf,
    UprightFont    = *-Roman,
    BoldFont       = *-Bold,
    ItalicFont     = *-Italic,
    BoldItalicFont = *-BoldItalic]
}{
  \IfFontExistsTF{Charter}{\setmainfont{Charter}}{
    \IfFontExistsTF{Iowan Old Style}{\setmainfont{Iowan Old Style}}{
      \IfFontExistsTF{Georgia}{\setmainfont{Georgia}}{}}}}

% Handwriting face, resolved once into \ppfont. Used for a heading or two (the
% answer key's "Answers"), not for the body and not for the drawings: the
% figures are read closely and want the text face.
%
% Naming a font that is not installed is a hard error, not a fallback, which is
% why this is a chain and not a line: seven sheets used to name one outright and
% stopped building the day it was uninstalled. Write \ppfont, never
% \fontspec{<a name>}. Excalifont is vendored in tools/fonts/ and installed by
% CI, so it is the one face that is always present.
\newcommand{\ppfont}{}
\IfFontExistsTF{Pretty Neat}{\renewcommand{\ppfont}{\fontspec{Pretty Neat}}}{
  \IfFontExistsTF{Humor Sans}{\renewcommand{\ppfont}{\fontspec{Humor Sans}}}{
    \IfFontExistsTF{Excalifont}{\renewcommand{\ppfont}{\fontspec{Excalifont}}}{
      \renewcommand{\ppfont}{}}}}

\setlength{\parindent}{0pt}
\setlength{\parskip}{0.38em}
\setlength{\emergencystretch}{1.2em}
\raggedbottom

% Break the page unless #1 of it is left.
%
% This is needspace's glue trick, written out because needspace.sty lives in
% texlive-latex-extra and not in a BasicTeX install. Do NOT rewrite it as a test
% on \pagegoal minus \pagetotal: the page builder has not necessarily run when
% the macro expands, so those numbers can still be the previous page's, the test
% says "no room" on a page that is empty, and the \newpage throws that page
% away. m02 shipped a blank page 3 that way, and it is invisible in the source.
\newcommand{\needroom}[1]{%
  \par
  \begingroup
    \setlength{\dimen0}{#1}%
    \vskip 0pt plus \dimen0
    \penalty -100%
    \vskip 0pt plus -\dimen0
  \endgroup}

% A part heading. The sheet has no title: it opens on Part 1. The heading
% carries the break test, because a heading printed at the foot of a page with
% nothing under it is worse than a short page.
\newcommand{\parthead}[2]{%
  \needroom{6\baselineskip}%
  \partheadhere{#1}{#2}}

% The same heading with no break test, for one already inside a minipage --
% \newpage means nothing in there.
\newcommand{\partheadhere}[2]{%
  \vspace{0.5em}
  {\large\bfseries Part #1 --- #2}\par
  \vspace{-0.4em}\rule{\textwidth}{0.8pt}\par
  \vspace{0.1em}}

% Centred fixed-width cell, for tables the student writes into.
\newcolumntype{C}[1]{>{\centering\arraybackslash}p{#1}}

% The lab the sheet hands off to, if it has one. Both the QR code and the
% printed address are the course's own short link, never the notebook's: paper
% outlives any one upload, so re-uploading a lab is a line on the server
% (/etc/caddy/Caddyfile, block go.skojaku.com) and not a reprint.
%
%   uvx --from segno segno --output=m07lab-qr.png --scale=20 --border=4 \
%       --error=m "https://go.skojaku.com/m07lab"
%
% --border=4 is the quiet zone the QR spec asks for and is not optional: at
% --border=1 some codes run to the edge of the PNG and OpenCV reads nothing.
\newcommand{\molablink}{https://go.skojaku.com/m07lab}
\newcommand{\molaburl}{\href{\molablink}{\texttt{go.skojaku.com/m07lab}}}
% The simulator, in the view the sheet needs (?sheet=1: the degree law hidden,
% the most and the least visited node named). Same Caddy block as the lab.
\newcommand{\simlink}{https://go.skojaku.com/m07sim}
\newcommand{\simurl}{\href{\simlink}{\texttt{go.skojaku.com/m07sim}}}

% xkcd-style hand-drawn line decoration.
% WARNING: applying [xkcd] to many nodes can throw "Dimension too large".
% Use it on a handful of \draw lines only, never on every node.
\makeatletter
\pgfset{
  /pgf/decoration/randomness/.initial=2,
  /pgf/decoration/wavelength/.initial=100
}
\pgfdeclaredecoration{sketch}{init}{
  \state{init}[width=0pt,next state=draw,persistent precomputation={
    \pgfmathsetmacro\pgf@lib@dec@sketch@t0
  }]{}
  \state{draw}[width=\pgfdecorationsegmentlength,
  auto corner on length=\pgfdecorationsegmentlength,
  persistent precomputation={
    \pgfmathsetmacro\pgf@lib@dec@sketch@t{mod(\pgf@lib@dec@sketch@t+pow(\pgfkeysvalueof{/pgf/decoration/randomness},rand),\pgfkeysvalueof{/pgf/decoration/wavelength})}
  }]{
    \pgfmathparse{sin(2*\pgf@lib@dec@sketch@t*pi/\pgfkeysvalueof{/pgf/decoration/wavelength} r)}
    \pgfpathlineto{\pgfqpoint{\pgfdecorationsegmentlength}{\pgfmathresult\pgfdecorationsegmentamplitude}}
  }
  \state{final}{}
}
\tikzset{xkcd/.style={decorate,decoration={sketch,segment length=0.5pt,amplitude=0.5pt}}}
\makeatother
\definecolor{pltrust}{HTML}{B14434}
\usetikzlibrary{calc}
% The three figures the Module 7 sheet draws, kept in their own file because
% the answer key draws exactly the same figures with the answers on them.
% Needs tikz and the colours pltblue and pltrust, which the including file sets
% up.
% ---------------------------------------------------------------------------
\tikzset{
  friend/.style  = {line width=1.3pt, black!58, line cap=round},
  person/.style  = {circle, draw=black!72, fill=white, line width=0.9pt,
                    minimum size=0.74cm, inner sep=0pt},
  % The walker's start: a second ring round the person, as on slide 6.
  start/.style   = {person, double=white, double distance=1.6pt},
  % A box the student writes a number into (a share, in Question 6).
  sbox/.style    = {draw=pltblue!70, line width=0.7pt, fill=white,
                    minimum width=1.25cm, minimum height=0.62cm, inner sep=0pt,
                    rounded corners=1.2pt},
  member/.style  = {circle, draw=black!72, fill=white, line width=0.9pt,
                    minimum size=0.66cm, inner sep=0pt, font=\small},
  dot/.style     = {circle, fill=pltblue!80!black, inner sep=0pt,
                    minimum size=0.22cm},
}

% ---------------------------------------------------------------------------
% Part 1: four friends. A--B, A--C, B--C, A--D. Degrees 3, 2, 2, 1, so 2m = 8
% and the long-run shares are 3/8, 2/8, 2/8, 1/8. D hangs off A, so the walker
% swings between A and the rest; the share column of Question 6 does not care.
%
% #1 is drawn last, on top: the answer key writes the shares into the boxes
% with \node at (bA) {...}; and so on.
\newcommand{\fournet}[1]{%
\begin{tikzpicture}
  \node[start]  (A) at (0, 0)      {A};
  \node[person] (B) at (1.9, 1.25) {B};
  \node[person] (C) at (2.6, -0.45){C};
  \node[person] (D) at (-1.9, 0)   {D};
  \draw[friend] (A) -- (B);
  \draw[friend] (A) -- (C);
  \draw[friend] (B) -- (C);
  \draw[friend] (A) -- (D);
  \node[font=\footnotesize, black!65, above=0.08cm of A] {start};
  \node[sbox] (bA) at (0, -1.05)   {};
  \node[sbox] (bB) at (3.35, 1.25) {};
  \node[sbox] (bC) at (2.6, -1.45) {};
  \node[sbox] (bD) at (-1.9, -1.0) {};
  #1
\end{tikzpicture}}

% ---------------------------------------------------------------------------
% Part 3: twelve members, the network of slides 14-45. Members 0-5 and 6-11
% are two groups of six, everybody friends with everybody in the group except
% 0-3 and 7-10, and two friendships cross: 2-8 and 5-11. 30 friendships.
% The layout is the deck's (slides/m07/figures/smallgraphs.py, TWO12_XY): two
% hexagons, the missing friendships are long diagonals, and the two crossing
% friendships are parallel horizontal lines that touch no other member.
%
% #1 is the style of members 0-5 and #2 of members 6-11 (empty on the sheet:
% the student darkens them; the key fills one side).
\newcommand{\tworooms}[2]{%
\begin{tikzpicture}[x=0.8cm, y=0.8cm]
  \foreach \n/\a in {2/60, 1/120, 0/180, 4/240, 5/300, 3/0}
    \coordinate (m\n) at ($(1.5, 2.0) + (\a:1.7)$);
  \foreach \n/\a in {6/60, 8/120, 10/180, 11/240, 9/300, 7/0}
    \coordinate (m\n) at ($(7.5, 2.0) + (\a:1.7)$);
  \foreach \i/\j in {0/1, 0/2, 0/4, 0/5, 1/2, 1/3, 1/4, 1/5, 2/3, 2/4, 2/5,
                     3/4, 3/5, 4/5,
                     6/7, 6/8, 6/9, 6/10, 6/11, 7/8, 7/9, 7/11, 8/9, 8/10,
                     8/11, 9/10, 9/11, 10/11,
                     2/8, 5/11}
    \draw[friend] (m\i) -- (m\j);
  \foreach \n in {0,...,5}  \node[member, #1] at (m\n) {\n};
  \foreach \n in {6,...,11} \node[member, #2] at (m\n) {\n};
\end{tikzpicture}}

% ---------------------------------------------------------------------------
% Part 3: the empty axes the second eigenvector is plotted on, one column per
% member, -1 to 1. #1 is drawn on top (the key's dots).
\newcommand{\eigaxes}[1]{%
\begin{tikzpicture}[x=1.15cm, y=2.9cm]
  \foreach \y in {-1, -0.5, 0.5, 1}
    \draw[black!14, line width=0.5pt] (-0.5, \y) -- (11.5, \y);
  \foreach \x in {0,...,11}
    \draw[black!14, line width=0.5pt] (\x, -1.12) -- (\x, 1.12);
  \draw[black!70, line width=0.8pt, dashed] (-0.5, 0) -- (11.5, 0);
  \draw[black!80, line width=0.9pt] (-0.5, -1.12) -- (-0.5, 1.12);
  \foreach \y/\l in {-1/{$-$1}, -0.5/{$-$0.5}, 0/{0}, 0.5/{0.5}, 1/{1}}
    \node[left, font=\small] at (-0.5, \y) {\l};
  \foreach \x in {0,...,11}
    \node[anchor=base, font=\small] at (\x, -1.27) {\x};
  \node[font=\small, black!65] at (5.5, -1.5) {member};
  #1
\end{tikzpicture}}


\begin{document}

\parthead{1}{A walker with no map}

\begin{minipage}[c]{0.5\textwidth}
\setlength{\parskip}{0.38em}
Four friends: A, B, C and D. A line joins two people who are friends.

A walker starts at A. At every step it picks one friend of the person it is
standing on, each friend with the same chance, and moves there. It has no map
and no memory of where it came from.
\end{minipage}
\hfill
\begin{minipage}[c]{0.46\textwidth}
\begin{center}
\resizebox{0.95\textwidth}{!}{\fournet{}}
\end{center}
\end{minipage}

\vspace{0.4em}

{\bf Question 1}: Before you compute anything, guess. After a very long walk,
at which of the four people is the walker found most often? Why?

{\bf Question 2}: In each cell of the table, write the chance that the walker
moves from the person of that row to the person of that column in one step.

\begin{center}
{\setlength{\tabcolsep}{6pt}
\begin{tabular}{|C{2.4cm}|C{1.7cm}|C{1.7cm}|C{1.7cm}|C{1.7cm}|}
\hline
\small from \textbackslash{} to & \bfseries A & \bfseries B & \bfseries C & \bfseries D \\
\hline
\bfseries A & & & & \\[0.3cm] \hline
\bfseries B & & & & \\[0.3cm] \hline
\bfseries C & & & & \\[0.3cm] \hline
\bfseries D & & & & \\[0.3cm] \hline
\end{tabular}}
\end{center}

Then add up each row. What do you get?

{\bf Question 3}: The walker starts at A. What is the chance that it is at B
after exactly three steps?

A route is the list of people the walker stands on, for example
A\,$\to$\,D\,$\to$\,A\,$\to$\,B. The chance of a route is the chances of its
three steps multiplied together: $\frac13 \times 1 \times \frac13 = \frac19$.

Write every three-step route from A to B in the table, with its chance. Then
add up the chances.

\begin{center}
{\setlength{\tabcolsep}{8pt}
\begin{tabular}{|C{5.0cm}|C{5.0cm}|}
\hline
\bfseries route & \bfseries chance \\ \hline
A\,$\to$\,D\,$\to$\,A\,$\to$\,B & $\frac13 \times 1 \times \frac13 = \frac19$ \\[0.08cm] \hline
& \\[0.3cm] \hline
& \\[0.3cm] \hline
& \\[0.3cm] \hline
& \\[0.3cm] \hline
\end{tabular}}
\end{center}

\parthead{2}{Drunk at the karate club}

Now a real network: a university karate club in the 1970s, with 34 members and
78 friendships between them. The same walker walks on it.

\begin{tcolorbox}[colback=white, boxsep=1pt, top=4pt, bottom=4pt]
\begin{minipage}[c]{0.24\textwidth}
\begin{center}
\includegraphics[width=2.7cm]{m07sim-qr.png}
\end{center}
\end{minipage}
\hfill
\begin{minipage}[c]{0.72\textwidth}
\textbf{Open the simulator}: point a camera at the square, or type \simurl.

On the screen, a member is a \textbf{node} and a friendship is an
\textbf{edge}.
\end{minipage}
\end{tcolorbox}

{\bf Question 4}: Press \textbf{Run 5,000 steps} four times: 20,000 steps.

Under the buttons, two boxes name the node with the most visits and the node
with the fewest. Each box gives the node's number of edges and its share of all
the visits. Write them in the table.

Then divide the edges by the share.

\begin{center}
{\setlength{\tabcolsep}{6pt}
\begin{tabular}{|C{3.2cm}|C{1.7cm}|C{1.7cm}|C{2.0cm}|C{2.6cm}|}
\hline
& \bfseries node & \bfseries edges & \bfseries share & \bfseries edges $\div$ share \\ \hline
\bfseries most visits & & & & \\[0.42cm] \hline
\bfseries fewest visits & & & & \\[0.42cm] \hline
\end{tabular}}
\end{center}

\needroom{5\baselineskip}%
{\bf Question 5}: Look at the two numbers in the last column of the table. Are
they the same? Compare them with your neighbour's.

The club has 78 friendships. How are your numbers related to 78?

\needroom{5\baselineskip}%
{\bf Question 6}: Go back to the four friends in Part 1. They have 4
friendships.

Use what you found in Question 5 to work out the share of time the walker
spends at A, B, C and D after a very long walk. Write the four shares in the
boxes on the drawing in Part 1.

Was your guess in Question 1 right?

\vfill
\begin{center}
\rule{\textwidth}{2pt}\\[0.3em]
{\large\bfseries STOP here. Wait for the lecture before you turn the page.}\\[0.1em]
\rule{\textwidth}{2pt}
\end{center}

\newpage

\parthead{3}{Two rooms}

Twelve members and 30 friendships. The walker is the same as in Part 1.

\begin{center}
\resizebox{0.92\textwidth}{!}{\tworooms{}{}}
\end{center}

\begin{tcolorbox}[colback=white, boxsep=1pt, top=4pt, bottom=4pt]
\begin{minipage}[c]{0.24\textwidth}
\begin{center}
\includegraphics[width=2.7cm]{m07lab-qr.png}
\end{center}
\end{minipage}
\hfill
\begin{minipage}[c]{0.72\textwidth}
\textbf{Open the notebook}: point a camera at the square, or type \molaburl.
A laptop or a phone will do; if you have neither, pair up.

It has one cell, printed below. Run it. The matrix \texttt{P} in it is the
table you wrote in Question 2, for these twelve members.
\end{minipage}
\end{tcolorbox}

\vspace{-0.3em}
{\small
\begin{verbatim}
import numpy as np

A = np.zeros((12, 12))                       # who is friends with whom
for group in [range(0, 6), range(6, 12)]:
    for i in group:
        for j in group:
            if i != j:
                A[i, j] = 1
for i, j in [(0, 3), (7, 10)]:
    A[i, j] = A[j, i] = 0
for i, j in [(2, 8), (5, 11)]:
    A[i, j] = A[j, i] = 1

P = A / A.sum(axis=1, keepdims=True)         # row i: the chances of leaving i
values, vectors = np.linalg.eig(P)
order = np.argsort(-values.real)             # largest eigenvalue first

print("eigenvalues:", np.round(values.real[order], 3) + 0.0)
v = vectors[:, order[1]].real                # the second eigenvector
print("second eigenvector:", np.round(v / np.abs(v).max(), 2) + 0.0)
\end{verbatim}
}

\needroom{14\baselineskip}%
{\bf Question 7}: The cell prints two lines: the twelve eigenvalues of
\texttt{P}, largest first, and the second eigenvector, one number per member.

Mark the twelve numbers of the second eigenvector as dots on the axes.

\begin{center}
\resizebox{0.9\textwidth}{!}{\eigaxes{}}
\end{center}

\needroom{4\baselineskip}%
{\bf Question 8}: On the drawing in Part 3, darken every member whose number is
above zero. What do the darkened members have in common?

\needroom{3\baselineskip}%
{\bf Question 9}: Which four members have the numbers closest to zero? Where
are they on the drawing?

\needroom{6\baselineskip}%
{\bf Question 10}: The second eigenvalue is 0.889. Work out
0.889\textsuperscript{5} and 0.889\textsuperscript{40}. A phone calculator is
fine.

Of the other eigenvalues, the one furthest from zero is $-$0.423. Work out
0.423\textsuperscript{5}.

After 5 steps, which eigenvalues still have a term in the walk? And after 40
steps?

\needroom{4\baselineskip}%
{\bf Question 11}: The lecture asked: short walks show communities, long walks
show degree. Why?

Answer in two sentences: one for the long walk, using Part 2, and one for the
short walk, using Part 3.

\needroom{6\baselineskip}%
{\bf Question 12} (finished early): Members 2 and 8 stop being friends. In the
code, delete \texttt{(2, 8),} from the last list and run the cell again.

Is the second eigenvalue now closer to 1, or further from 1?

Does a walker that starts at member 0 stay among members 0 to 5 for longer, or
for less long? Why?

\end{document}
